% a_entries.tex - Terms beginning with A
\chapter*{A}
\addcontentsline{toc}{chapter}{A}

\dictentry{Account Number}{/əˈkaʊnt ˈnʌmbər/, \textit{n.}}{
A unique sequence of numbers assigned to a cardholder account, identifying
the issuing bank or financial institution and the type of financial-
transaction card.
}
\indexentry{Account Number}

\dictentry{Accounts Receivable Conversion (ARC)}{/əˈkaʊnts rɪˈsiːvəbəl kənˈvɝːʒən/, \textit{n.}}{
A process by which paper checks submitted for payment are electronically
converted into data for processing through the Automated Clearing House
(ACH).
}
\indexentry{Accounts Receivable Conversion (ARC)}

\dictentry{Acculynk}{/ˈækjuːlɪŋk/, \textit{n.}}{
A suite of payment solutions by First Data that includes True Debit\texttrademark{}
gateway, PaySecure\textregistered{}, and Payzur, enabling secure and
cost-effective eCommerce transactions and peer-to-peer disbursements via
Electronic Funds Transfer (EFT) networks.
}
\indexentry{Acculynk}

\dictentry{ACH}{/eɪ siː eɪtʃ/, \textit{abbr.}}{
Automated Clearing House. An electronic network for financial transactions in the United States. ACH processes large volumes of credit and debit transactions in batches, making it cost-effective for recurring payments like payroll, direct deposits, and bill payments.

\example{Your employer uses ACH to deposit your salary directly into your bank account every two weeks.}

\crossref{\textbf{Direct Deposit}, \textbf{Wire Transfer}, \textbf{Batch Processing}}

\etymology{The ACH network was established in the 1970s to create a more efficient alternative to paper checks.}
}
\indexentry{ACH}

\dictentry{ACH Rules}{/eɪ siː eɪtʃ ruːlz/, \textit{n.}}{
The NACHA Operating Rules and Operating Guidelines governing the interregional
exchange and settlement of ACH transactions.
}
\indexentry{ACH Rules}

\dictentry{Acquirer/Acquiring Bank}{/əˈkwaɪər əˈkwaɪərɪŋ bæŋk/, \textit{n.}}{
A financial institution that processes payment transactions, such as credit or
debit card payments, on behalf of a merchant, underwriting the merchant
account and providing necessary hardware and software.
}
\indexentry{Acquirer/Acquiring Bank}

\dictentry{Acquiring}{/əˈkwaɪərɪŋ/, \textit{n.}}{
The process by which a bank or financial institution (the acquiring bank)
enables a merchant to accept card payments for goods or services.
}
\indexentry{Acquiring}

\dictentry{Address Verification Service (AVS)}{/əˈdrɛs ˌvɛrɪfɪˈkeɪʃən sɝːvɪs/, \textit{n.}}{
A security system allowing merchants to verify a cardholder's billing
address with the issuing bank during card-not-present transactions, such
as online purchases, to reduce fraud.
}
\indexentry{Address Verification Service (AVS)}

\dictentry{Adjustments}{/əˈdʒʌstmənts/, \textit{n.}}{
Corrections made by the acquirer or processor to address discrepancies,
such as disputes or duplicate charges, in payment transactions.
}
\indexentry{Adjustments}

\dictentry{Affinity Card}{/əˈfɪnəti kɑːrd/, \textit{n.}}{
A credit card issued in conjunction with an organization, such as a
professional group or alumni association, where the issuer pays a royalty
to the organization for marketing the card to its members.
}
\indexentry{Affinity Card}

\dictentry{Alipay\texttrademark{}}{/ˈælɪpeɪ/, \textit{n.}}{
A mobile and online payment platform developed by Alibaba Group\textregistered{},
facilitating payment transactions both online and in stores.
}
\indexentry{Alipay}

\dictentry{Alternative Payments/Local Payments}{/ɔːlˈtɝːnətɪv ˈpeɪmənts ˈloʊkəl ˈpeɪmənts/, \textit{n.}}{
Payment types that bypass traditional card associations, offering
alternatives to credit card transactions, such as debit cards, prepaid
cards, bank transfers, and mobile payments.
}
\indexentry{Alternative Payments/Local Payments}

\dictentry{American Express\textregistered{} (AMEX)}{/əˈmɛrɪkən ɛkˈsprɛs/, \textit{n.}}{
A multinational financial services corporation headquartered in New York
City, offering payment, travel, and expense-management solutions for
individuals and businesses.
}
\indexentry{American Express (AMEX)}

\dictentry{APDU}{/ˌeɪ piː diː ˈjuː/, \textit{abbr.}}{
Application Protocol Data Unit. The communication unit between a smart card
reader and a smart card. It consists of a command sent by the reader and a
response sent by the card. This "question and answer" format is the fundamental
language of chip-based payments.

\example{When a payment terminal asks a card for its list of supported
applications, it sends an APDU command, and the card sends back an APDU
response containing the list.}

\analogy{Biological}{Like a \textbf{Neurotransmitter} crossing a synapse. One
neuron (the reader) sends a specific chemical packet (the command) to another
neuron (the card), which then reacts and sends a signal back. This precise
exchange of messages is how the two "cells" of the payment system communicate.}

\crossref{\textbf{ISO/IEC 7816}, \textbf{Smart Card}, \textbf{OpenCard Framework (OCF)}}

\etymology{Defined in the ISO/IEC 7816-4 standard for integrated circuit cards
with contacts.}
}
\indexentry{APDU}

\dictentry{Arbitration}{/ˌɑːrbɪˈtreɪʃən/, \textit{n.}}{
A process used by credit card companies to determine whether the acquirer
(or entity processing payment transactions for a merchant) is responsible
for a chargeback.
}
\indexentry{Arbitration}

\dictentry{ATM Interchange Fee}{/eɪ tiː ɛm ˈɪntərtʃeɪndʒ fiː/, \textit{n.}}{
A fee paid by the issuer member to the acquirer member for processing an
ATM transaction through the STAR\textregistered{} Network or similar
networks.
}
\indexentry{ATM Interchange Fee}

\dictentry{ATM Terminal Driving}{/eɪ tiː ɛm ˈtɝːmɪnəl ˈdraɪvɪŋ/, \textit{n.}}{
Maintaining the connection and software between an Automated Teller
Machine (ATM) and relevant parties, ensuring proper transaction
processing.
}
\indexentry{ATM Terminal Driving}

\dictentry{Automated Teller Machine (ATM)}{/ˈɔːtəmeɪtɪd ˈtɛlər məˈʃiːn/, \textit{n.}}{
An electronic banking terminal that allows customers to complete basic
transactions, such as cash withdrawals, without the presence of a bank
representative.
}
\indexentry{Automated Teller Machine (ATM)}

\dictentry{Authentication}{/ɔːˌθɛntɪˈkeɪʃən/, \textit{n.}}{
The process of verifying the identity of a credit card user, typically
through a Personal Identification Number (PIN) or other secure means.
}
\indexentry{Authentication}

\dictentry{Authorization}{/ˌɔːθərəˈzeɪʃən/, \textit{n.}}{
The initial request by a merchant for a customer's card issuer to release
funds, verifying that the card has sufficient funds to cover the
transaction.
}
\indexentry{Authorization}

\dictentry{Authorization Approval Code}{/ˌɔːθərəˈzeɪʃən əˈpruːvəl koʊd/, \textit{n.}}{
A code indicating that a transaction is approved and the card may be
honored by the merchant.
}
\indexentry{Authorization Approval Code}

\dictentry{Authorization Code}{/ˌɔːθərəˈzeɪʃən koʊd/, \textit{n.}}{
A unique identifier provided by the issuing bank, confirming the approval
of a transaction during the authorization process.
}
\indexentry{Authorization Code}

\dictentry{Authorization Response}{/ˌɔːθərəˈzeɪʃən rɪˈspɑːns/, \textit{n.}}{
The response from the issuer indicating whether a transaction is approved,
declined, or requires additional verification.
}
\indexentry{Authorization Response}

% Additional A terms will be added here
